\section*{الفصل \ref{ch:b1:layered-earth} --- الأرض ذات الطبقات}

\begin{solution}{exo:b1:layered-earth:1}
نسب الحجم هي مكعبات نسب أنصاف الأقطار: اللب
$(3480/6371)^3 = 0.1630$، أي \qty{16.3}{\percent}؛ \omterm{def:b1:layered-earth:layers}{اللب الداخلي}
$(1221.5/6371)^3 = 0.00705$، \qty{0.70}{\percent}.
\end{solution}

\begin{solution}{exo:b1:layered-earth:2}
فوق: $\Phi = 13.7166^2 - \tfrac43\times7.26466^2 = 188.145 - 70.367 =
\qty{117.78}{km^2.s^{-2}}$، $K = \rho\Phi = \qty{5566.45}{kg.m^{-3}}\times
\qty{117.78e6}{m^2.s^{-2}} = \qty{655.6}{GPa}$. وتحت: $\Phi = 8.06482^2 =
\qty{65.04}{km^2.s^{-2}}$، $K = 9903.49\times\num{65.04e6} = \qty{644.1}{GPa}$.
يكاد معامل الحجم لا يتغيّر (\qty{-1.7}{\percent}). ما يتغيّر هو معامل القصّ، $\mu = \rho v_S^2 = \qty{293.8}{GPa}$ فوقه و$0$ تحته (لأنه سائل)، والكثافة، التي ترتفع بنسبة \qty{78}{\percent}: ومن هنا هبوط
$v_P$ رغم تساوي $K$ تقريبًا.
\end{solution}

\begin{solution}{exo:b1:layered-earth:3}
$\sin i_c = 6.2/8.1 = 0.7654$، $i_c = \ang{49.95}$. المسافة الحرجة
$2\times30\times\tan\ang{49.95} = 60\times1.1906 = \qty{71.4}{km}$. \omterm{def:b1:layered-earth:head-wave}{مسافة التقاطع}
$60\sqrt{14.3/1.9} = 60\times2.743 = \qty{164.6}{km}$.
\end{solution}

\begin{solution}{exo:b1:layered-earth:4}
$g = Gm/r^2 = \num{6.674e-11}\times\num{1.94e24}/(\num{3.48e6})^2 =
\qty{10.69}{m.s^{-2}}$، أي أكبر ب\qty{9}{\percent} من السطح. وعند النزول إلى اللب يهبط مربع نصف القطر إلى $(3480/6371)^2 = 0.298$ من قيمته عند السطح، بينما تهبط الكتلة الجاذبة إلى $1.94/5.972 = 0.325$ من المجموع فقط، لأن اللب يحوي ثلث الكتلة: $0.325/0.298 = 1.09$.
\end{solution}

\begin{solution}{exo:b1:layered-earth:5}
قرب المصدر $v_1 = 60/10.0 = \qty{6.0}{km.s^{-1}}$. وبعيدًا
$v_2 = (300 - 200)/(45.8 - 33.3) = 100/12.5 = \qty{8.0}{km.s^{-1}}$؛ وزمن القطع
$t_i = 33.3 - 200/8.0 = \qty{8.3}{s}$. ثم $\sin i_c =
0.75$، و$\cos i_c = 0.6614$، و$h = t_i v_1/(2\cos i_c) = 8.3\times6.0/1.3229
= \qty{37.6}{km}$. تحقّق: $x_c = 2\times37.6\times\sqrt{14/2} = \qty{199}{km}$،
حيث تصل الموجتان معًا عند $199/6.0 = 199/8.0 + 8.3 = \qty{33.2}{s}$.
\end{solution}

\begin{solution}{exo:b1:layered-earth:6}
لكثافة منتظمة $\rho$، $I = \tfrac{8\pi}{3}\rho\int_0^R r^4\,\mathrm dr =
\tfrac{8\pi}{15}\rho R^5$ و$M = \tfrac43\pi\rho R^3$، إذن
$I/MR^2 = (8/15)/(4/3) = 0.4$. طبقتان ($r_c = \qty{3.48e6}{m}$،
$R = \qty{6.371e6}{m}$): $M = \tfrac43\pi\bigl(11000\,r_c^3 + 4450\,(R^3 -
r_c^3)\bigr) = \qty{5.977e24}{kg}$، أعلى ب\qty{0.08}{\percent} من الكتلة المقيسة؛ و$I = \tfrac{8\pi}{15}\bigl(11000\,r_c^5 + 4450\,(R^5 - r_c^5)\bigr)$ يعطي
$I/MR^2 = 0.3457$. فالكتلة صحيحة، لكن المعامل أعلى ب\qty{4.5}{\percent} مما ينبغي: \omterm{def:b1:layered-earth:layers}{فالوشاح} الحقيقي ليس منتظمًا، وكثافته تزداد مع العمق، فيدفع الكتلة نحو الداخل.
\end{solution}

\begin{solution}{exo:b1:layered-earth:7}
$\Phi = 12.16^2 - \tfrac43\times6.673^2 = 147.87 - 59.37 =
\qty{88.49}{km^2.s^{-2}}$. وبعمق $z$،
$\mathrm d\rho/\mathrm dz = \rho g/\Phi = 4860\times9.9/\num{88.49e6} =
\qty{5.44e-4}{kg.m^{-4}}$، أي \qty{0.544}{kg.m^{-3}} لكل كيلومتر، أو
\qty{0.54}{g.cm^{-3}} لكل \qty{1000}{km}. ترتفع PREM بمقدار $5.357 - 4.844 =
0.513$ عبر هذه \qty{1000}{km}، أي أقل ب\qty{6}{\percent}: فالتدرج المحسوب عند القمة يُطبَّق على الفترة كلها، بينما يزداد $\Phi$ مع العمق ويتناقص التدرج.
\end{solution}

\begin{solution}{exo:b1:layered-earth:8}
$P \approx \bar\rho g z = 3570\times9.9\times\num{670e3} = \qty{23.7}{GPa}$
(في PREM: \qty{23.9}{GPa}). إنه الضغط، نحو \qty{24}{GPa}، الذي ينقسم عنده الرينغوودايت إلى البريدجمانيت والفيروبيريكلاس: انفصال \qty{660}{km}.
\end{solution}

\begin{solution}{exo:b1:layered-earth:9}
(أ)~الحديد السائل: $10.90\times1.05 = 11.4$ إلى $10.90\times1.10 =
\qty{12.0}{g.cm^{-3}}$. (ب)~$12.764/12.166 = 1.049$: قفزة قدرها
\qty{4.9}{\percent} في PREM، وهي أكبر قليلًا من \qty{4.5}{\percent} في الدراسات الحديثة. (ج)~التجمّد وحده يجعل الصلب أكثف من السائل بقدر ضئيل؛ والقفزة الأكبر تعني أن الصلب يحوي عناصر خفيفة أقل من السائل: فهي تفضّل السائل، والتجمّد عند \omterm{def:b1:layered-earth:discontinuity}{حدّ اللب الداخلي} يترك معظمها خلفه.
\end{solution}

\begin{solution}{exo:b1:layered-earth:10}
لنضع $x_* = 2h\tan i_c$. فإن $\sqrt{x_*^2 + 4h^2} = 2h/\cos i_c$، و
$t_r(x_*) = 2h/(v_1\cos i_c)$. وتعطي \omterm{def:b1:layered-earth:head-wave}{موجة الرأس}
$t(x_*) = 2h\tan i_c/v_2 + 2h\cos i_c/v_1 = \frac{2h}{v_1}\bigl(\sin^2 i_c/\cos
i_c + \cos i_c\bigr) = 2h/(v_1\cos i_c)$، باستعمال $\sin i_c = v_1/v_2$: النتيجة نفسها.
ميل القطع الزائد هو $t_r'(x) = x/(v_1\sqrt{x^2 + 4h^2})$؛ وعند $x_*$ يساوي
$\sin i_c/v_1 = 1/v_2$، ميل خط \omterm{def:b1:layered-earth:head-wave}{موجة الرأس}. وبما أن
$t_r'' > 0$، فإن القطع الزائد محدّب ويقع فوق كل مماس له: لأجل
$x > x_*$، $t_r > t_{\text{الرأس}}$. وعند المسافة الحرجة يصطدم الانعكاس بالسطح الفاصل بالزاوية $i_c$ تمامًا، وهي نشأة \omterm{def:b1:layered-earth:head-wave}{موجة الرأس}؛ وبعدها تسبق \omterm{def:b1:layered-earth:head-wave}{موجة الرأس}، التي تكسب بسيرها في الوسط الأسرع، الموجة الأخرى.
\end{solution}

\begin{solution}{exo:b1:layered-earth:11}
$m(r) = \tfrac43\pi\rho r^3$، إذن $g = Gm/r^2 = \tfrac43\pi G\rho r$. ثم
$P(r) = \int_r^R \rho\,\tfrac43\pi G\rho s\,\mathrm ds = \tfrac23\pi G\rho^2(R^2
- r^2)$. عند المركز $P = \tfrac23\pi\times\num{6.674e-11}\times5513^2\times
(\num{6.371e6})^2 = \qty{172}{GPa}$، أي \qty{47}{\percent} من \qty{363.85}{GPa} في PREM: فتركيز الكتلة نحو المركز يضاعف الضغط المركزي تقريبًا.
\end{solution}

\begin{solution}{exo:b1:layered-earth:12}
مع $z$ العمق، $\mathrm d\rho/\mathrm dz = \rho g/\Phi$؛ فإذا كان $g/\Phi$ ثابتًا $a$، فإن $\mathrm d(\ln\rho)/\mathrm dz = a$ و$\rho = \rho_0
e^{a(z - z_0)}$. هنا $a(z - z_0) = 7.9\times\num{2258.5e3}/\num{87e6} =
0.2051$ و$\rho = 9.90\times e^{0.2051} = 9.90\times1.2276 =
\qty{12.15}{g.cm^{-3}}$، أي أدنى ب\qty{0.2}{\percent} من $12.17$ في PREM: فكل الزيادة في الكثافة عبر \omterm{def:b1:layered-earth:layers}{اللب الخارجي} انضغاط ذاتي.
\end{solution}

\begin{solution}{pb:b1:layered-earth:1}
\textbf{1.} $V = \tfrac43\pi(\num{6.371e6})^3 = \qty{1.0832e21}{m^3}$.

\textbf{2.} $\tfrac43\pi(r_2^3 - r_1^3)$: \omterm{def:b1:layered-earth:layers}{القشرة} \omterm{def:b1:layered-earth:layers}{والوشاح} العلوي ومنطقة الانتقال
\qty{3.0706e20}{m^3}؛ \omterm{def:b1:layered-earth:layers}{الوشاح} السفلي \qty{5.9961e20}{m^3}؛ \omterm{def:b1:layered-earth:layers}{اللب الخارجي}
\qty{1.6890e20}{m^3}؛ \omterm{def:b1:layered-earth:layers}{اللب الداخلي} \qty{7.634e18}{m^3}.

\textbf{3.} \qty{28.35}{\percent}، \qty{55.36}{\percent}، \qty{15.59}{\percent}
و\qty{0.70}{\percent}: \omterm{def:b1:layered-earth:layers}{الوشاح} السفلي هو أكبر القشور، أكثر من نصف
حجم الأرض.

\textbf{4.} $15.59 + 0.70 = \qty{16.3}{\percent}$.

\textbf{5.} \omterm{def:b1:layered-earth:layers}{القشرة} \omterm{def:b1:layered-earth:layers}{والوشاح} العلوي ومنطقة الانتقال $\num{3.0706e20}\times3570 =
\qty{1.0962e24}{kg}$؛ \omterm{def:b1:layered-earth:layers}{الوشاح} السفلي \qty{2.9381e24}{kg}؛ \omterm{def:b1:layered-earth:layers}{اللب الخارجي}
\qty{1.8410e24}{kg}؛ \omterm{def:b1:layered-earth:layers}{اللب الداخلي} \qty{9.841e22}{kg}.

\textbf{6.} \qty{5.9737e24}{kg}، أعلى ب\qty{0.03}{\percent} من الكتلة المقيسة
\qty{5.972e24}{kg}.

\textbf{7.} $\num{5.9737e24}/\num{1.0832e21} = \qty{5515}{kg.m^{-3}}$.

\textbf{8.} $\num{9.841e22}/\num{5.9737e24} = \qty{1.65}{\percent}$.

\textbf{9.} $(1.0962 + 2.9381)/5.9737 = \qty{67.5}{\percent}$.

\textbf{10.} اللب هو \qty{16.3}{\percent} من الحجم و
\qty{32.5}{\percent} من الكتلة، أي ضعف ذلك: كثافته المتوسطة،
\qty{10986}{kg.m^{-3}}، تبلغ نحو $2.5$ ضعف كثافة \omterm{def:b1:layered-earth:layers}{الوشاح}.

\textbf{11.} \omterm{def:b1:layered-earth:layers}{قشرة} ذات كثافة منتظمة تعطي
$I = \tfrac{8\pi}{3}\rho\int_{r_1}^{r_2} r^4\,\mathrm dr = \tfrac{8\pi}{15}\rho
(r_2^5 - r_1^5)$.

\textbf{12.} \omterm{def:b1:layered-earth:layers}{القشرة} \omterm{def:b1:layered-earth:layers}{والوشاح} العلوي ومنطقة الانتقال \qty{2.676e37}{kg.m^2}؛ \omterm{def:b1:layered-earth:layers}{الوشاح} السفلي
\qty{4.525e37}{kg.m^2}؛ \omterm{def:b1:layered-earth:layers}{اللب الخارجي} \qty{9.27e36}{kg.m^2}؛ \omterm{def:b1:layered-earth:layers}{اللب الداخلي}
\qty{5.87e34}{kg.m^2}.

\textbf{13.} \omterm{def:b1:layered-earth:layers}{الوشاح} السفلي: له نصف الكتلة تقريبًا، عند أنصاف أقطار كبيرة. \omterm{def:b1:layered-earth:layers}{واللب الداخلي}
يحوي \qty{1.6}{\percent} من الكتلة، وكلها قرب المحور: حصته \qty{0.07}{\percent}.

\textbf{14.} $I = \qty{8.1345e37}{kg.m^2}$ و$I/MR^2 = \num{8.1345e37}/
(\num{5.9737e24}\times(\num{6.371e6})^2) = 0.3355$.

\textbf{15.} أعلى ب\qty{1.4}{\percent} من $0.3308$، وأدنى كثيرًا من $0.4$: فالنموذج يقترب من التركيز الصحيح للكتلة نحو المركز.

\textbf{16.} نعم. داخل كل \omterm{def:b1:layered-earth:layers}{قشرة} تكون الكثافة الحقيقية أعلى عند القاع منها عند القمة؛ \omterm{def:b1:layered-earth:layers}{فالقشرة} المنتظمة ذات الكتلة نفسها تضع كثيرًا من تلك الكتلة قرب نصف قطرها الخارجي، حيث تُحسب ب$r^2$، فتبالغ في تقدير $I$.

\textbf{17.} $g = GM/R^2 = \num{6.674e-11}\times\num{5.9737e24}/
(\num{6.371e6})^2 = \qty{9.822}{m.s^{-2}}$.

\textbf{18.} كتلة اللب هي $\num{1.8410e24} + \num{9.841e22} =
\qty{1.9394e24}{kg}$، و$g = G m/r_c^2 = \qty{10.69}{m.s^{-2}}$.

\textbf{19.} $(\num{1.0962e24} + \num{2.9381e24})/(\num{3.0706e20} +
\num{5.9961e20}) = \num{4.0343e24}/\num{9.0667e20} = \qty{4450}{kg.m^{-3}}$.

\textbf{20.} $P \approx 4450\times10\times\num{2.891e6} = \qty{128.6}{GPa}$،
أدنى ب\qty{5}{\percent} من \qty{135.75}{GPa} في PREM.

\textbf{21.} $\rho_c = \num{1.9394e24}/\num{1.7653e20} = \qty{10986}{kg.m^{-3}}$،
و$\Delta P = \tfrac23\pi\times\num{6.674e-11}\times10986^2\times
(\num{3.48e6})^2 = \qty{204.3}{GPa}$.

\textbf{22.} $128.6 + 204.3 = \qty{332.9}{GPa}$، أدنى ب\qty{8.5}{\percent} من
\qty{363.85}{GPa} في PREM. وكلا التقديرين يهملان زيادة الكثافة مع العمق داخل كل منطقة، وهي التي تضع كتلة أكبر في الأعماق، وتُبقي الجاذبية أعلى في النزول، وترفع $\rho g$، أي تدرّج الضغط.

\textbf{23.} الكتلة الكلية (السؤال 6)، ومعامل عزم القصور الذاتي
(السؤال 15) والجاذبية عند السطح (السؤال 17، مقابل متوسط \qty{9.82}{m.s^{-2}} المقيس) تُقارن بقياسات الأرض؛ أما الجاذبية عند اللب والضغوط فلا يمكن مقارنتها إلا بPREM، وهو نموذج بذاته.

\textbf{24.} يحوي اللب \textbf{\qty[text-series-to-math=true]{32.5}{\percent}}
من كتلة الأرض.
\end{solution}
